\chapter*{Аннотация}
% \addcontentsline{toc}{chapter}{Аннотация}

В работе исследована интерференционная картина колец Ньютона как метод для определения радиуса искривления линзы и разности длин волн двух спектральных компонент света на примере ртутной лампы. Получены значения радиусов колец в зависимости от их порядкового номера, благодаря чему измерен радиус линзы. Определено расстояние между центрами двух четких систем колец, полученных в результате ''биения'' двух цветовых компонент, что позволило вычислить разницу длин волн зеленого и желтого цветов света ртутной лампы. Полученное значение совпало в пределах погрешности с табличным. Результаты подтвердили применимость данного метода для бесконтактного определения радиусов линз и разности длин волн двух спектральных компонент исследуемого света.
